\section{What makes a shared world testable}\label{sec:open}

The construction separates four questions. Which observables are common is
answered by the intersection of admissible hulls. Which state is common is
answered by restrictions of the actual process state, or by a compatibility
condition when source states are independent. Which records agree is a joint
instrument question. Which reports are right depends on the standards and
references that connect those records to the target. None of the four can
replace the others.

The same-source route advances Foundation's implementation responsibility.
It derives a largest common algebra relative to given operations, provides
a Type~III model with no foreground overlap, and connects the algebra to a
noise-only persistence mechanism. It does not explain why those operations,
seeds or permissions first form. Foundation~III treats the return of bounded
feedback to the source; T-DOME treats retained histories and self-calibration;
Brain Anchor treats learned sensorimotor relations. These mechanisms must
change an actual common process before they can be used to explain the origin
of the sharing structure rather than its compatibility.

The reference construction supplies a further link. Every non-scalar common
$\Gamma_3$-hull has a candidate finite positive-energy excitation. With legal
controls and paid readings it can calibrate energy orientation; with the
additional higher-dimensional source and response conditions it can test the
geometric law. This joins Foundation~I's selection, Foundation~II's sharing
and Consciousness~II's identification. The null translation generator and
physical energy must not be silently equated. Nor does common stress determine
a unique metric or common experience. HAFF-D's sourced response and the
Q-RAIF control interface require their own physical couplings and costs.

The target-specific criterion strengthens this connection. Public correctness
needs enough reference separation for the public target, rather than identical
private frames or deletion of old histories. A shared decoded calibration
can support different subsequent decisions because each reader also retains
different task information. The same energy reading can be public even when
private charge orientations remain undetermined.

We propose that the operational extent of a shared world is selected partly
by the references its participants can jointly sustain. This is a mechanism
claim: source changes affect the availability and price of reference
preparations; changed reference risk alters the legal common task family;
the resulting actions change the next source. It predicts loss of some common
capabilities before every stored private history disappears. The finite
reference witness establishes the calibration and preservation part, and
Foundation~III supplies finite causal return. A general derivation still
requires the same device to supply the reference, actual source and subsequent
readout at a feasible maintained budget. Uncalibrated common symmetry errors
and the zero-gain fault are direct failure conditions. A strict hull inclusion
alone supplies no universal positive heat price; joint output entropy,
reusable memory and the actual heat channel control that price.

Several sharp boundaries remain. Nontriviality of an arbitrary shared hull
is not guaranteed; the two-register example makes it fail. A common error
survives whenever the whole paid protocol retains the relevant symmetry.
Energy references cannot remove internal transformations commuting with
energy. Finite sample and finite time change accessibility continuously even
when an exact fixed algebra jumps. The comparison needed next is therefore
between task risks and legal resource vectors, not just between exact
intersections. The calibration tables already show why consensus is an
inadequate proxy: agreement can improve while an agreeing pair becomes less
reliable, and a channel that passes useful information can also pass a fault.

For independent sources the compatibility requirements remain in the
appendices. A unique state-preserving expectation follows from modular
invariance, but a common projected memory kernel or a canonical private
complement does not. Establishing those objects from interactions, and
estimating the reference or repair policy from the readers' own use history,
are concrete further formation problems. They preserve the original question
of how a shared world and its readers arise together.

